% Policy on Prevention of Sexual Harassment at the Workplace (POSH).
%
% NOT derived from a .docx -- there was no POSH policy in the set. Added
% because the Sexual Harassment of Women at Workplace (Prevention,
% Prohibition and Redressal) Act, 2013 applies to every workplace in
% India regardless of headcount, and INTERNS are expressly within the
% definition of an aggrieved person at a workplace.
%
% Where fewer than ten workers are employed, no Internal Committee is
% required and complaints go to the District Local Committee. That is
% what registry/org.yaml `posh.internal_committee: false` selects below.
%
% WORDING IS A DRAFT FOR REVIEW BY QUALIFIED COUNSEL. This document
% provides the structure and the routing; it is not a legal opinion.
\documentclass{paridyum}
\begin{document}

\paridyumletterhead
\paridyumdateref{\dkDocDate}{\dkDocRef}
\paridyumtitle{Policy on Prevention of Sexual Harassment (POSH)}

\textit{This policy applies to everyone who works at or with Paridyum, including
interns, trainees, contractors, consultants and volunteers, whether paid or unpaid,
and whether working on-site, remotely or in a hybrid arrangement.}

\section*{1. Commitment}
Paridyum is committed to providing a workplace free from sexual harassment. Sexual
harassment is a serious form of misconduct. Any person who works at or with Paridyum
is entitled to be treated with dignity, and to raise a concern without fear of
disadvantage.

\section*{2. Scope}
This policy applies at Paridyum's premises and at any place visited in the course of
work, including transport arranged for work, off-site meetings, work events, and
online spaces used for work such as chat, video calls and shared documents. Remote
work does not place conduct outside this policy.

\section*{3. What Sexual Harassment Means}
Sexual harassment includes any unwelcome act or behaviour of a sexual nature, whether
direct or by implication, such as:
\begin{itemize}[leftmargin=*, noitemsep]
    \item physical contact or advances;
    \item a demand or request for sexual favours;
    \item making sexually coloured remarks;
    \item showing pornography; and
    \item any other unwelcome physical, verbal or non-verbal conduct of a sexual
          nature, including through messages, email, calls or images.
\end{itemize}

The following, where they occur in relation to sexual conduct, also amount to sexual
harassment:
\begin{itemize}[leftmargin=*, noitemsep]
    \item an implied or explicit promise of preferential treatment;
    \item an implied or explicit threat of detrimental treatment;
    \item an implied or explicit threat about present or future status;
    \item interference with work, or creating an intimidating, offensive or hostile
          environment; and
    \item humiliating treatment likely to affect health or safety.
\end{itemize}

Whether conduct is unwelcome is judged from the point of view of the person
experiencing it, not from the intention of the person doing it.

\section*{4. How to Raise a Complaint}
\ifdkOrgPoshInternalCommittee
A complaint may be made in writing to Paridyum's Internal Committee. Details of the
Internal Committee are displayed at the workplace and are available from
\dkOrgPoshContactEmail.
\else
Paridyum currently employs fewer than ten people and therefore does not have an
Internal Committee. \textbf{A complaint of sexual harassment may be made directly to
the Local Committee constituted by the District Officer:}

\vspace{0.5em}
\begin{infobox}
\textbf{Local Committee}\\[0.4em]
\dkOrgPoshLocalCommittee

\vspace{0.5em}
A complaint may also be raised internally, in confidence, by writing to
\textbf{\dkOrgPoshContactEmail}. Raising it internally does not remove the right to
go to the Local Committee, and Paridyum will assist with making that complaint if
asked.
\end{infobox}
\fi

\vspace{0.5em}
A complaint should ordinarily be made within three months of the incident, or within
three months of the last incident in a series. This period may be extended where there
were circumstances that prevented an earlier complaint.

\section*{5. Support and Assistance}
Paridyum will provide reasonable assistance to any person who wishes to make a
complaint, including help with putting the complaint in writing and with contacting
the Local Committee. Paridyum will assist a complainant who wishes to make a police
complaint, and will not obstruct any person from doing so.

\section*{6. Confidentiality}
The identity of the complainant, the respondent, the witnesses, and the content of a
complaint and its outcome will be kept confidential and will not be published or made
known to the public, press or media.

\section*{7. No Retaliation}
No person will be disadvantaged for making a complaint in good faith, for supporting a
complaint, or for giving evidence. Retaliation is itself misconduct and will be acted
on. This applies fully to interns: an internship will not be shortened, altered, or
ended, and a certificate or reference will not be withheld, because a person raised a
concern.

\section*{8. Action on Findings}
Where a complaint is upheld, Paridyum will act on the recommendations made, which may
include ending the internship or engagement of the person responsible, and any other
action available under applicable law.

\section*{9. Awareness}
This policy is provided to every intern at the start of the internship, and is
acknowledged in writing as part of the onboarding documentation.

\section*{10. Acknowledgement}
By signing below, the Intern confirms that they have received and read this policy and
know how to raise a complaint under it.

\vspace{1.5em}
\paridyumofficers{\dkSigCeoName}{\dkSigCeoDesignation}{\dkSigCtoName}{\dkSigCtoDesignation}

\vspace{1.5em}
\paridyumcountersign{Intern}{\dkPIdentityLegalName}

\end{document}
